\section{Scene Representation Taxonomy}
\label{sec:representation}

The choice of scene representation is arguably the single most consequential
design decision in any \nvs{} system.
It determines the fundamental trade-offs a method faces---rendering speed,
reconstruction accuracy, memory footprint, generalisability, and
editability---and has a direct bearing on whether the representation can
serve as a persistent world state for downstream planning or action conditioning.
Over the past several years, the field has explored a rich variety of
representations, each occupying a different position on the
speed--compactness--quality Pareto frontier.
We identify four canonical families below and return to each in depth in
Sections~\ref{sec:nerf}--\ref{sec:generalisable}.

\textbf{Implicit representations} (MLPs, hash grids) encode the scene as a
continuous function $f_\theta(\mathbf{x})$ evaluated at arbitrary 3-D query
points.
Their key advantages are compactness---a complete scene is stored as a
few megabytes of network weights---and continuous differentiability, which
supports high-quality reconstruction from dense observations and enables
gradients to flow through the rendering process for downstream optimisation.
Their central limitation is rendering speed: each pixel requires integrating
along a ray through many sequential network evaluations, making real-time
rendering impractical without architectural acceleration.
NeRF~\citep{mildenhall2020nerf} and its successors (Section~\ref{sec:nerf})
constitute the canonical family, with subsequent work substituting the MLP
for multi-resolution hash grids~\citep{muller2022instant} or tensor
decompositions~\citep{chen2022tensorf} to substantially accelerate training.

\textbf{Explicit representations} (voxels, point clouds, Gaussians) store scene
attributes at discrete spatial elements, enabling tile-based or splat-based
rasterisation that achieves real-time frame rates through massive GPU
parallelism.
The trade-off is memory: a faithful 3DGS scene may occupy hundreds of
megabytes to gigabytes, and the discrete nature of the representation
makes it less amenable to continuous differentiable reasoning.
3D Gaussian Splatting~\citep{kerbl2023gaussian} (Section~\ref{sec:3dgs}) is
the dominant current representative, having largely supplanted point-cloud-based
methods by incorporating learnable appearance and anisotropic geometry.

\textbf{Hybrid representations} combine an explicit spatial structure---typically
a triplane~\citep{chan2022efficient} or multi-resolution
grid~\citep{muller2022instant}---with a shallow MLP decoder.
They strike a practical balance between the quality of implicit models and
the speed of explicit ones: the spatial structure provides memory-efficient
feature caching, while the MLP decoder retains the capacity for high-frequency
appearance modelling.
This combination has proven especially powerful for large-scale generalised
reconstruction, and it is the architectural backbone of most Large Reconstruction
Models~(Section~\ref{sec:lrm}) and EG3D-style 3-D-aware
GANs~\citep{chan2022efficient}.

\textbf{Compositional representations} decompose the scene into semantic
entities (objects, backgrounds, agents) with independent sub-representations.
This structural inductive bias enables object-level editing and compositional
generation~\citep{niemeyer2021giraffe,zhou2024gala3d}, but typically requires
supervision or segmentation to separate entities.
It is also the natural representation for world models that must reason about
independently acting objects---vehicles, humans, and manipulable items---each
with its own dynamics model.
The growing intersection of compositional representations and world modeling
is one of the most active research fronts we identify in Section~\ref{sec:open}.

\medskip
Each of the four sections that follow (Sections~\ref{sec:nerf}--\ref{sec:generalisable})
develops one branch of this taxonomy in detail.
The taxonomy is not exhaustive, and the boundaries between families are blurring:
feed-forward Gaussian methods~(Section~\ref{sec:generalisable}) simultaneously
exhibit the speed of explicit representations and the generalisability of
implicit ones, while SDF-based neural methods~(Section~\ref{sec:nerf}) occupy
an intermediate position between implicit and explicit through their mesh
extraction pipeline.
We return to the implications of this convergence in Section~\ref{sec:world}.

% ─────────────────────────────────────────────────────────────────────────────
